% GENERATED by tools/gen.py from data/cards/grammar.tsv -- do not edit.
\section{A1}
\Card{A1}{Wymowa}{1}{Jak czytasz „ei” w słowie „mein”?}{jak „aj”: [majn]}
\Card{A1}{Wymowa}{1}{Jak czytasz „ie” w słowie „sieben”?}{jak długie „i”: [zibn]}
\Card{A1}{Wymowa}{1}{Jak czytasz „eu” w słowie „Deutsch”?}{jak „oj”: [dojcz]}
\Card{A1}{Wymowa}{1}{Jak czytasz „sch” i „sp-/st-” na początku słowa?}{„sch” = „sz”; „sp-” = „szp-”, „st-” = „szt-” (Schule, Sprache, Stadt)}
\Card{A1}{Wymowa}{2}{Kiedy „ch” jest miękkie, a kiedy twarde?}{Twarde po a, o, u, au (Buch, Nacht); miękkie po pozostałych samogłoskach i po l, n, r (ich, Milch)}
\Card{A1}{Wymowa}{1}{Jak czytasz „z” i „v”?}{„z” = „c” (zehn); „v” zwykle = „f” (Vater)}
\Card{A1}{Wymowa}{2}{Dlaczego „Tag” czytamy [tak]?}{Na końcu słowa b, d, g tracą dźwięczność: [p], [t], [k]}
\Card{A1}{Powitania}{1}{Jak powiesz „Dzień dobry” formalnie po południu?}{Guten Tag}
\Card{A1}{Powitania}{1}{Jak zapytasz „Jak się Pan/Pani miewa?”}{Wie geht es Ihnen?}
\Card{A1}{Powitania}{2}{Kiedy używasz „Sie”, a kiedy „du”?}{„Sie” formalnie (klient, urząd, rozmowa kwalifikacyjna); „du” w zespole i wśród znajomych, po propozycji}
\Card{A1}{Czasowniki}{1}{Odmień „sein” w liczbie pojedynczej.}{ich bin, du bist, er/sie/es ist}
\Card{A1}{Czasowniki}{1}{Odmień „haben” w liczbie pojedynczej.}{ich habe, du hast, er/sie/es hat}
\Card{A1}{Czasowniki}{1}{Jakie końcówki mają regularne czasowniki w Präsens?}{-e, -st, -t, -en, -t, -en}
\Card{A1}{Czasowniki}{2}{Podaj „du” od „nehmen”, „fahren”, „lesen”.}{du nimmst, du fährst, du liest}
\Card{A1}{Szyk}{1}{Na którym miejscu stoi odmieniony czasownik w zdaniu oznajmującym?}{Na drugim: Heute arbeite ich zu Hause.}
\Card{A1}{Szyk}{2}{Co to jest klamra zdaniowa?}{Odmieniony czasownik na 2. miejscu, bezokolicznik lub przedrostek na końcu: Ich muss früh aufstehen.}
\Card{A1}{Rodzajniki}{1}{Jakie rodzajniki określone mają trzy rodzaje i liczba mnoga?}{der (m), die (f), das (n), die (mnoga)}
\Card{A1}{Rodzajniki}{2}{Co zmienia się w bierniku (Akkusativ)?}{Tylko rodzaj męski: der $\rightarrow$ den, ein $\rightarrow$ einen, kein $\rightarrow$ keinen}
\Card{A1}{Rodzajniki}{2}{Który rodzaj mają rzeczowniki na -ung, -heit, -keit, -schaft, -ion?}{żeński (die)}
\Card{A1}{Przeczenie}{2}{Kiedy „kein”, a kiedy „nicht”?}{„kein” przeczy rzeczownik z „ein”/bez rodzajnika; „nicht” czasownik, przymiotnik, rzeczownik z „der/die/das”}
\Card{A1}{Zwroty}{1}{Jak powiesz „Mam 30 lat”?}{Ich bin 30 Jahre alt. (sein, nie haben)}
\Card{A1}{Zwroty}{1}{Jak grzecznie zamówisz kawę?}{Ich möchte einen Kaffee. / Ich hätte gern einen Kaffee.}
\Card{A1}{Czas}{1}{Co znaczy „halb neun”?}{8:30 (pół do dziewiątej)}
\Card{A1}{Czas}{2}{Które przyimki stoją przy dniu, godzinie i miesiącu?}{am Montag, um acht Uhr, im Mai}
\Card{A1}{Modalne}{2}{Czym różni się „Du musst nicht kommen” od „Du darfst nicht kommen”?}{Pierwsze: nie musisz; drugie: nie wolno ci}
\Card{A1}{Przyimki}{3}{Kiedy po „in/auf/an” jest celownik, a kiedy biernik?}{Wo? (miejsce) $\rightarrow$ Dativ; Wohin? (ruch do celu) $\rightarrow$ Akkusativ}
\Card{A1}{Czasowniki}{2}{Gdzie stoi przedrostek „auf” w „aufstehen”?}{Na końcu zdania: Ich stehe um sieben Uhr auf.}
\Card{A1}{Imperativ}{2}{Jak utworzysz rozkaz do „du” i „Sie” od „kommen”?}{Komm! / Kommen Sie!}
\Card{A1}{Liczby}{2}{Jak powiesz 21, 34, 67?}{einundzwanzig, vierunddreißig, siebenundsechzig (jedności przed dziesiątkami)}
\Card{A1}{Szwajcaria}{1}{Jak pozdrowisz w Szwajcarii i jak podziękujesz potocznie?}{Grüezi; Merci (dialektalne, ale zrozumiałe)}
\Card{A1}{Szwajcaria}{1}{Czy w Szwajcarii używa się „ß”?}{Nie: zawsze „ss” (Strasse, dreissig)}
\section{A2}
\Card{A2}{Przypadki}{2}{Które przyimki zawsze wymagają celownika?}{aus, bei, mit, nach, seit, von, zu, gegenüber}
\Card{A2}{Przypadki}{2}{Które przyimki zawsze wymagają biernika?}{durch, für, gegen, ohne, um}
\Card{A2}{Przypadki}{2}{Wymień zaimki osobowe w celowniku.}{mir, dir, ihm, ihr, ihm, uns, euch, ihnen / Ihnen}
\Card{A2}{Przypadki}{3}{Jaka jest kolejność dopełnień w zdaniu?}{Dwa rzeczowniki: Dativ przed Akkusativ; dwa zaimki: Akkusativ przed Dativ (Ich gebe ihn ihm)}
\Card{A2}{Perfekt}{2}{Z jakim czasownikiem pomocniczym tworzysz Perfekt czasowników ruchu?}{sein: Ich bin nach Zürich gefahren.}
\Card{A2}{Perfekt}{2}{Jak tworzysz Partizip II regularnego czasownika?}{ge- + rdzeń + -t: gemacht, gearbeitet (bez ge- przy -ieren i be-/ver-/er-/ent-)}
\Card{A2}{Perfekt}{2}{Podaj Partizip II: gehen, kommen, sehen, essen.}{gegangen, gekommen, gesehen, gegessen}
\Card{A2}{Präteritum}{2}{Które czasowniki mówisz zwykle w Präteritum?}{sein (war), haben (hatte) i modalne (konnte, musste)}
\Card{A2}{Nebensätze}{2}{Gdzie stoi odmieniony czasownik po „weil”, „dass”, „wenn”, „ob”?}{Na końcu zdania podrzędnego}
\Card{A2}{Nebensätze}{3}{Czym różni się „weil” od „denn”?}{„weil” $\rightarrow$ czasownik na końcu; „denn” $\rightarrow$ szyk normalny}
\Card{A2}{Stopniowanie}{2}{Podaj stopniowanie: gut, viel, gern.}{besser/am besten; mehr/am meisten; lieber/am liebsten}
\Card{A2}{Stopniowanie}{2}{„als” czy „wie” przy porównaniu?}{Wyższy stopień: als (größer als); równość: so … wie}
\Card{A2}{Zwrotne}{2}{Podaj zaimki zwrotne dla wszystkich osób.}{mich, dich, sich, uns, euch, sich}
\Card{A2}{Zwrotne}{3}{Jaki przyimek po „sich freuen” (przyszłość) i „sich interessieren”?}{sich freuen auf + Akk; sich interessieren für + Akk}
\Card{A2}{Przymiotniki}{3}{Jakie końcówki przymiotnika po „der/die/das” w mianowniku?}{-e (m, f, n); w mnogiej -en}
\Card{A2}{Przymiotniki}{3}{Jakie końcówki po „ein” w mianowniku męskim i nijakim?}{-er (ein neuer Laptop), -es (ein neues Auto)}
\Card{A2}{Futur}{2}{Jak tworzysz Futur I?}{werden + bezokolicznik na końcu: Ich werde Deutsch lernen.}
\Card{A2}{Uprzejmość}{2}{Podaj trzy grzeczne formy próśb.}{Könnten Sie …? / Ich hätte gern … / Wäre es möglich, …?}
\Card{A2}{E-mail}{2}{Jak zaczniesz formalny e-mail do pani Keller?}{Sehr geehrte Frau Keller,}
\Card{A2}{Szwajcaria}{2}{Jak kończysz formalny e-mail w Szwajcarii?}{Freundliche Grüsse (DE: Mit freundlichen Grüßen)}
\Card{A2}{Szwajcaria}{1}{Jak po szwajcarsku: rower, bilet, parkować?}{Velo, Billett, parkieren}
\section{B1}
\Card{B1}{Nebensätze}{3}{„damit” czy „um … zu”?}{Ten sam podmiot: um … zu; różne podmioty: damit}
\Card{B1}{Nebensätze}{3}{„als” czy „wenn” w przeszłości?}{als: jednorazowe zdarzenie; wenn: powtarzalne}
\Card{B1}{Nebensätze}{2}{Jaki czas po „nachdem”?}{Perfekt (przy Präsens w głównym) lub Plusquamperfekt (przy Präteritum)}
\Card{B1}{Relativsätze}{3}{Skąd bierzesz rodzaj i przypadek zaimka względnego?}{Rodzaj i liczbę z poprzednika, przypadek z roli w zdaniu względnym}
\Card{B1}{Relativsätze}{3}{Podaj zaimki względne w celowniku.}{dem (m, n), der (f), denen (mnoga)}
\Card{B1}{Genitiv}{2}{Jakie rodzajniki w dopełniaczu?}{des (m, n) + -s, der (f, mnoga)}
\Card{B1}{Genitiv}{2}{Podaj przyimki z dopełniaczem.}{wegen, trotz, statt, während, innerhalb, außerhalb}
\Card{B1}{Adjektive}{3}{Jak odmieniasz przymiotnik bez rodzajnika?}{Końcówka rodzajnika: frischer Kaffee, frische Milch, frisches Brot}
\Card{B1}{Passiv}{2}{Jak tworzysz Passiv Präsens i Präteritum?}{werden/wurde + Partizip II: Der Code wird getestet. / wurde getestet.}
\Card{B1}{Passiv}{3}{Jak dodajesz wykonawcę w Passiv?}{von + Dativ (osoba), durch + Akkusativ (środek, przyczyna)}
\Card{B1}{Konjunktiv II}{2}{Podaj formy „wäre”, „hätte”, „würde”, „könnte”.}{ich wäre, hätte, würde, könnte (er tak samo)}
\Card{B1}{Konjunktiv II}{3}{Jak zbudujesz nierealny warunek w teraźniejszości?}{Wenn ich Zeit hätte, würde ich kommen. (jedna „würde”)}
\Card{B1}{Konjunktiv II}{3}{Jak w przeszłości nierealnej?}{hätte/wäre + Partizip II: Wenn ich das gewusst hätte, wäre ich gekommen.}
\Card{B1}{Präteritum}{2}{Podaj Präteritum: gehen, kommen, sehen, finden.}{ging, kam, sah, fand}
\Card{B1}{Plusquamperfekt}{3}{Jak tworzysz Plusquamperfekt?}{hatte/war + Partizip II}
\Card{B1}{Präpositionen}{3}{Jaki przyimek: warten …, denken …, sprechen …, fragen …?}{warten auf (Akk), denken an (Akk), sprechen über (Akk), fragen nach (Dat)}
\Card{B1}{Präpositionen}{3}{Jak tworzysz pytanie o rzecz i odpowiedź zaimkiem?}{wo(r)- + przyimek: Worauf wartest du? / da(r)- + przyimek: Ich warte darauf.}
\Card{B1}{Konnektoren}{3}{Jaki szyk po „deshalb”, „trotzdem”, „außerdem”?}{Czasownik na 2. miejscu: Deshalb haben wir …}
\Card{B1}{Meinung}{2}{Jak wyrazisz opinię i sprzeciw uprzejmie?}{Ich finde, dass … / Da bin ich anderer Meinung. / Ich sehe das etwas anders.}
\Card{B1}{Bewerbung}{2}{Jakie trzy pytania z rozmowy kwalifikacyjnej musisz umieć?}{Erzählen Sie etwas über sich. / Warum möchten Sie bei uns arbeiten? / Was sind Ihre Stärken?}
\Card{B1}{IT}{2}{Podaj Partizip II: mergen, deployen, reviewen.}{gemergt, deployt, reviewt}
\Card{B1}{IT}{2}{Jak powiesz w standupie, że jesteś zablokowany?}{Ich bin blockiert, weil mir der Zugriff fehlt.}
\Card{B1}{Szwajcaria}{2}{Czym jest Franchise i Prämie w ubezpieczeniu?}{Franchise: udział własny rocznie; Prämie: składka (im wyższa Franchise, tym niższa Prämie)}
\Card{B1}{Szwajcaria}{2}{Co znaczy Znüni, Zmittag, Zvieri?}{przekąska ok. 9:00, obiad, przekąska po południu}
\Card{B1}{Szwajcaria}{2}{Czy Szwajcarzy mówią dialektem do osoby uczącej się?}{Zwykle przechodzą na standard; dialekt rozumiej, standardem mów}
\section{B2}
\Card{B2}{Nominalstil}{3}{Zamień „Wir prüfen den Antrag” na styl nominalny.}{Die Prüfung des Antrags erfolgt (durch uns).}
\Card{B2}{Nominalstil}{3}{Podaj Funktionsverbgefüge dla „entscheiden” i „bereitstellen”.}{eine Entscheidung treffen; zur Verfügung stellen}
\Card{B2}{Partizipien}{3}{Co to jest Partizip I i II jako określenie?}{I: trwa (die steigenden Kosten); II: zakończone/bierne (die gespeicherten Daten)}
\Card{B2}{Partizipien}{3}{Zamień na określenie: die Daten, die im Test verwendet wurden.}{die im Test verwendeten Daten}
\Card{B2}{Konjunktiv I}{3}{Kiedy używasz Konjunktiv I?}{W mowie zależnej (dystans, prasa): Er sagte, er sei krank.}
\Card{B2}{Konjunktiv I}{3}{Co robisz, gdy Konjunktiv I = Indikativ?}{Zastępujesz Konjunktiv II: Sie sagten, sie hätten keine Zeit.}
\Card{B2}{Passiv}{3}{Czym różni się „Die Tür wird geschlossen” od „Die Tür ist geschlossen”?}{proces (Vorgangspassiv) vs stan (Zustandspassiv)}
\Card{B2}{Passiv}{3}{Podaj trzy zastępniki strony biernej.}{man; sich lassen (Das lässt sich lösen); -bar; sein + zu (ist zu beheben)}
\Card{B2}{Konnektoren}{3}{Jak działa „je … desto”?}{Je + zdanie podrzędne (czasownik na końcu), desto + zdanie główne (Je mehr …, desto besser wird …)}
\Card{B2}{Konnektoren}{3}{Podaj spójniki: sposób, skutek, warunek.}{indem; sodass; falls / sofern / es sei denn}
\Card{B2}{Konnektoren}{2}{Zamień „ohne dass er fragt” na bezokolicznik.}{ohne zu fragen}
\Card{B2}{Präsentation}{2}{Jak odwołasz się do slajdu i podsumujesz?}{Wie Sie auf dieser Folie sehen, … / Zusammenfassend lässt sich sagen, dass …}
\Card{B2}{Präsentation}{3}{„um” czy „auf” przy wzroście i spadku?}{um 20 Prozent gestiegen (o ile); von 80 auf 50 gesunken (do ilu)}
\Card{B2}{Vertrag}{2}{Zapytaj, czy wynagrodzenie jest brutto czy netto.}{Ist das Gehalt brutto oder netto?}
\Card{B2}{Vertrag}{3}{Jak uprzejmie sygnalizujesz niezadowolenie z oferty?}{Das entspricht nicht ganz meinen Vorstellungen.}
\Card{B2}{Dokumentation}{2}{Podaj nagłówki raportu o incydencie.}{Zusammenfassung, Zeitleiste, Grundursache, Maßnahmen, Nächste Schritte}
\Card{B2}{Gespräch}{3}{Jak zaczniesz trudną rozmowę bez oskarżeń?}{Darf ich kurz etwas ansprechen? Mir ist aufgefallen, dass …}
\Card{B2}{Feinheiten}{3}{Jaką funkcję mają partykuły „doch”, „mal”, „eben/halt”?}{doch: przypomnienie oczywistości; mal: łagodzi prośbę; eben/halt: „tak już jest”}
\Card{B2}{Feinheiten}{2}{Jak czytasz długi wyraz złożony?}{Od końca: ostatni człon wyznacza rodzaj i znaczenie podstawowe}
\Card{B2}{Szwajcaria}{3}{Czym jest Quellensteuer i kiedy dotyczy?}{Podatek u źródła; dotyczy często pracowników zagranicznych, zależy od pozwolenia, kantonu i dochodu (sprawdź u źródła)}
\Card{B2}{Szwajcaria}{3}{Jak zachowuje się Szwajcar wobec krytyki w pracy?}{Rzeczowo, uprzejmie, w cztery oczy; unika ostrych sformułowań}
\Card{B2}{Egzamin}{2}{Jakie cztery moduły ma egzamin B2?}{Lesen, Hören, Schreiben, Sprechen}
\Card{B2}{Egzamin}{2}{Jak postępujesz z nieznanym słowem w nagraniu?}{Nie zatrzymuję się, słucham dalej i wnioskuję z kontekstu}
